\section{The CrossLeak Attack}
\label{sec_CrossLeak}


 \begin{figure}
 	\centering
 	\includegraphics[width=0.999\linewidth]{../../../../../绘图/2026/INFOCOM27/Figure_2}
 	\caption{Overview of the CrossLeak attack. CrossLeak compares the original and unlearned aggregated models on a public probe set, trains a sparse crosscoder on paired representations, and infers deleted labels and deletion-sensitive features from the crosscoder generated statistical differences.}
 	\label{fig:figure2}
 \end{figure}
 


\subsection{Overview}
\label{sec:approach_overview}
CrossLeak is a post-unlearning privacy attack that infers what has been removed from a federated model. The key observation is that unlearning does not merely change the final prediction scores. It also changes the internal representation geometry: directions that were strongly supported by the deleted data are weakened, displaced, or made less class-consistent after unlearning. These changes can remain visible even when the unlearning protocol hides client updates through secure aggregation and even when the released model preserves normal test accuracy.



CrossLeak has two preparation stages. First, it collects paired pre- and post-unlearning representations from the two models on the same probe inputs. Second, it trains a sparse CrossCoder to decompose the paired representations into shared features and unlearning-specific changes. Then, CrossLeak scores labels by aggregating representation-level and prediction-level evidence. CrossLeak can also recover deletion-sensitive features and maps them to human-interpretable concepts using top-activating examples and CLIP-based concept matching. We demonstrate an overview process of CrossLeak in \Cref{fig:figure2}.


\subsection{Paired Model Signals}
\label{sec:approach_paired_signals}

 
For each probe input $x_i\in D_p$, CrossLeak extracts both logits and intermediate representations:
\begin{equation}
	g_i^o = f_{\theta_o}(x_i), \quad g_i^u = f_{\theta_u}(x_i),
	\label{eq:paired_logits}
\end{equation}
\begin{equation}
	z_i^o = h_{\theta_o}(x_i), \quad z_i^u = h_{\theta_u}(x_i),
	\label{eq:paired_representations}
\end{equation}
Here $h_{\theta}(\cdot)$ is the representation layer used by the attack. CrossLeak trains the CrossCoder on paired representations and keeps logits as prediction-level evidence for label inference. CrossLeak constructs a paired representation vector
\begin{equation}
    q_i = [\bar z_i^o; \bar z_i^u],
    \label{eq:crossleak_input}
\end{equation}
where $\bar z_i^o$ and $\bar z_i^u$ are normalized representations. The paired form is important: using only $f_{\theta_u}$ reveals what the final model knows, whereas CrossLeak needs the direction and magnitude of what changed.


\subsection{Sparse CrossCoder for Unlearning Differences}
\label{sec:approach_crosscoder}

The central component of CrossLeak is a sparse CrossCoder trained on paired representations. The CrossCoder learns a dictionary of latent features that jointly reconstructs the before- and after-unlearning signals while separating shared information from unlearning-sensitive information. Given $q_i$, the encoder produces a sparse code $a_i\in \mathbb{R}^{m}$:
\begin{equation}
    a_i = \sigma(W_e q_i + b_e),
    \label{eq:crosscoder_encoding}
\end{equation}
where $W_e$ and $b_e$ are the encoder weight matrix and bias, and $\sigma(\cdot)$ is a nonnegative activation. The decoder reconstructs the two representation views:
\begin{equation}
    [\hat z_i^o;\hat z_i^u] = W_d a_i + b_d.
    \label{eq:crosscoder_decoding}
\end{equation}
Here $W_d$ and $b_d$ are the decoder weight matrix and bias. We view $W_d$ as two decoder blocks, $W_d^o$ and $W_d^u$, which reconstruct the before- and after-unlearning representations, respectively.

CrossLeak trains the CrossCoder with reconstruction, sparsity, and change-awareness objectives:
\begin{equation}
\begin{aligned}
    \mathcal{L}_{cc}
    =& \frac{1}{n}\sum_{i=1}^{n}
    \left(
        \|\hat z_i^o-\bar z_i^o\|_2^2
        + \|\hat z_i^u-\bar z_i^u\|_2^2
    \right)  \\
    &+ \lambda_1 \frac{1}{n}\sum_{i=1}^{n}\|a_i\|_1
    + \lambda_{\Delta}\frac{1}{n}\sum_{i=1}^{n}
        \|\widehat{\Delta z_i} - \Delta z_i\|_2^2,
\end{aligned}
\label{eq:crosscoder_loss}
\end{equation}
where $\Delta z_i=\bar z_i^o-\bar z_i^u$ and $\widehat{\Delta z_i}=\hat z_i^o-\hat z_i^u$. The first term makes the learned features faithful to the model representations, the second term encourages a compact set of interpretable features, and the third term forces the dictionary to preserve the representation displacement induced by unlearning.

After training, each feature $j$ has a before-side decoder vector $w_j^o$ and an after-side decoder vector $w_j^u$. CrossLeak measures the unlearning sensitivity of feature $j$ by
\begin{equation}
    s_j^{\Delta} = \|w_j^o-w_j^u\|_2 \cdot
    \frac{1}{n}\sum_{i=1}^{n} a_{ij}.
    \label{eq:feature_delta_score}
\end{equation}
Features with large $s_j^{\Delta}$ are active on probe samples and have a strong before-after discrepancy, making them candidates for deletion-sensitive concepts.



\subsection{Deleted Label Inference}
\label{sec:approach_label_inference}

CrossLeak infers deleted labels by combining two families of signals: representation displacement and confidence degradation. For each class $c$, let $D_p^c$ be the probe samples labeled as class $c$. The representation displacement score is
\begin{equation}
    S_{\Delta z}(c)=
    \frac{1}{|D_p^c|}\sum_{x_i\in D_p^c}
    \|\bar z_i^o-\bar z_i^u\|_2.
    \label{eq:delta_z_score}
\end{equation}
This score is high when the class manifold moves substantially after unlearning.

The confidence-drop score measures whether the original model's support for class $c$ is selectively reduced:
\begin{equation}
    S_{\mathrm{conf}}(c)=
    \frac{1}{|D_p^c|}\sum_{x_i\in D_p^c}
    \left(p_{i,c}^{o}-p_{i,c}^{u}\right).
    \label{eq:confidence_drop_score}
\end{equation}
where $p_{i,c}^{o} = \operatorname{softmax}(f_{\theta_o}(x_{i,c}))$ is the prediction probability of samples on class $c$. This score captures deletion evidence that appears at the output layer. It is especially useful when unlearning directly suppresses the target class prediction without causing a large global representation shift.

 
The attacker ranks labels using one of these heads or a normalized combination:
\begin{equation}
    S_{\mathrm{label}}(c)=
    \alpha S_{\Delta z}^{*}(c)
    + \beta S_{\mathrm{conf}}^{*}(c),
    \label{eq:label_score}
\end{equation}
where $S^{*}$ denotes min-max normalization across labels. If the deletion budget $k$ is known, CrossLeak outputs the top-$k$ labels. Otherwise, it uses a score gap or threshold rule to determine the predicted deletion set.

\subsection{Deletion-Sensitive Feature Inference}
\label{sec:approach_feature_inference}

Beyond labels, CrossLeak identifies which latent features were most affected by unlearning. Here, a feature denotes a sparse crosscoder feature, i.e., a learned dictionary direction in the paired representation space, which may correspond to a visual attribute, texture, class-specific pattern, or trigger-like cue. A feature is considered deletion-sensitive if it is (i) active on probe samples, (ii) substantially changed between the before- and after-unlearning decoders, and (iii) mainly associated with samples from the inferred deleted labels. CrossLeak combines these requirements through a feature ranking score:
\begin{equation}
    S_{\mathrm{feat}}(j)=
    s_j^{\Delta}\cdot
    R_j^{\mathrm{before}}\cdot
    C_j^{\mathrm{target}},
    \label{eq:feature_score}
\end{equation}
where
\begin{equation}
    R_j^{\mathrm{before}}=
    \frac{\|w_j^o\|_2}{\|w_j^o\|_2+\|w_j^u\|_2+\epsilon}
    \label{eq:before_strength}
\end{equation}
is the before-exclusive strength ratio. Here $\epsilon$ is a small positive constant used only for numerical stability when both decoder norms are close to zero. A larger $R_j^{\mathrm{before}}$ means that feature $j$ has more mass in the before-unlearning decoder than in the after-unlearning decoder, which is consistent with deleted information being suppressed by unlearning. The target-label concentration term is
\begin{equation}
    C_j^{\mathrm{target}} =
    \frac{1}{|T_j|}
    \sum_{x_i\in T_j}
    \mathbb{I}[\hat y_i\in \hat Y_u]
    \label{eq:target_coverage}
\end{equation}
where $T_j$ is the set of top-activating probe samples for feature $j$, $\hat Y_u$ is the inferred deleted-label set, and $\hat y_i=\arg\max f_{\theta_o}(x_i)$ is the label assigned by the original model. The final sensitive feature set $\hat F_u$ is obtained by ranking $S_{\mathrm{feat}}(j)$ and selecting the top features.

 


\noindent
\textbf{Semantic Concept Recovery.} Latent features are informative but not immediately human-readable. CrossLeak interprets a sensitive feature $j$ by collecting its top activating probe samples $T_j$ and comparing them with a vocabulary of candidate concepts. For image tasks, CrossLeak uses a CLIP encoder to embed both the top activating images and textual prompts. Let $\phi_I(\cdot)$ be the CLIP image encoder and $\phi_T(\cdot)$ be the CLIP text encoder. For concept $r$, the semantic score is
\begin{equation}
    S_{\mathrm{clip}}(j,r)=
    \frac{1}{|T_j|}\sum_{x_i\in T_j}
    \cos\left(\phi_I(x_i),\phi_T(r)\right).
    \label{eq:clip_concept_score}
\end{equation}
The highest scoring concepts are used as explanations of the recovered feature. 

%For backdoor unlearning, the same mechanism can reveal the trigger-like visual pattern because trigger-sensitive features tend to have high before-exclusive strength and coherent top activating samples.


\begin{algorithm}[h]
	\caption{CrossLeak Attack} \label{alg:crossleak}
	\begin{normalsize}
		\BlankLine
		\KwIn{Original model $f_{\theta_o}$; unlearned model $f_{\theta_u}$; probe set $D_p$; prediction budget $k$; concept vocabulary $\mathcal{V}$.}
		\KwOut{Inferred deleted labels $\hat Y_u$; deletion-sensitive features $\hat F_u$; semantic concepts $\hat C_u$.}
		
		\BlankLine
		\tcp{(A) Collect paired model evidence}
		\ForEach{$x_i\in D_p$}{
			$g_i^o \gets f_{\theta_o}(x_i)$, \quad $g_i^u \gets f_{\theta_u}(x_i)$\;
			$z_i^o \gets h_{\theta_o}(x_i)$, \quad $z_i^u \gets h_{\theta_u}(x_i)$\;
			$p_i^o \gets \operatorname{softmax}(g_i^o)$, \quad $p_i^u \gets \operatorname{softmax}(g_i^u)$\;
			$q_i \gets [\bar z_i^o;\bar z_i^u]$\;
		}
		
		\BlankLine
		\tcp{(B) Train sparse CrossCoder on paired representations}
		Train the CrossCoder by minimizing $\mathcal{L}_{cc}$ in \Cref{eq:crosscoder_loss}\;
		\ForEach{CrossCoder feature $j$}{
			Compute displacement score $s_j^\Delta$ using \Cref{eq:feature_delta_score}\;
			Collect top-activating probe samples $T_j$\;
		}
		
		\BlankLine
		\tcp{(C) Deleted-label inference}
		\ForEach{candidate label $c\in\mathcal{Y}$}{
			Compute $S_{\Delta z}(c)$ using \Cref{eq:delta_z_score}\;
			Compute $S_{\mathrm{conf}}(c)$ using \Cref{eq:confidence_drop_score}\;
			Compute $S_{\mathrm{label}}(c)$ using \Cref{eq:label_score}\;
		}
		$\hat Y_u \gets \operatorname{TopK}_{c\in\mathcal{Y}} S_{\mathrm{label}}(c)$ with budget $k$\;
		
		\BlankLine
		\tcp{(D) Deletion-sensitive feature inference}
		\ForEach{CrossCoder feature $j$}{
			Compute $R_j^{\mathrm{before}}$ and $C_j^{\mathrm{target}}$ using \Cref{eq:before_strength,eq:target_coverage}\;
			Compute $S_{\mathrm{feat}}(j)$ using \Cref{eq:feature_score}\;
		}
		$\hat F_u \gets \operatorname{TopK}_{j} S_{\mathrm{feat}}(j)$\;
		
		\ForEach{feature $j\in \hat F_u$}{
			Match $T_j$ to concepts in $\mathcal{V}$ using $S_{\mathrm{clip}}$ in \Cref{eq:clip_concept_score}\;
		}
		$\hat C_u \gets$ top matched concepts for $\hat F_u$\;
		
		\BlankLine
		\Return $\hat Y_u$, $\hat F_u$, and $\hat C_u$\;
	\end{normalsize}
\end{algorithm}


\subsection{CrossLeak Attack Pipeline}
\label{sec:approach_pipeline}

The complete CrossLeak pipeline is summarized in \Cref{alg:crossleak}. The attack takes the original model $f_{\theta_o}$, the unlearned model $f_{\theta_u}$, and a public probe set $D_p$ as input. It first constructs paired model evidence by querying both models on the same probe samples. The paired representations are used to train the sparse CrossCoder, while logits and probabilities are retained for prediction-level label inference.

After the CrossCoder is trained, CrossLeak performs two inference tasks. For deleted-label inference, it ranks labels using representation displacement and confidence drop, and outputs the top-$k$ labels when the deletion budget is known. For deletion-sensitive feature inference, it ranks CrossCoder features according to their before-after displacement, before-exclusive strength, and target-label concentration. Finally, CrossLeak maps the top-ranked features to semantic concepts using CLIP-based concept matching, producing both the inferred deleted labels and human-interpretable feature evidence.


